\subsection{The relative action class of the controlled-$Z$ state}

\begin{lemma}[Coefficient maps and local reduction identities]
    \label{lem:mpug_reduction_coefficient_map}
    \lean{MPSTensor.complexOfRing_kronecker,
        MPSTensor.isReduction_complexOfRing_of_local_compression}
    \leanok
    Let $f:R\to\C$ be a homomorphism from a commutative ring. It preserves
    matrix Kronecker products. Suppose matrices over $R$ satisfy
    \[
        VW=I,\qquad VB^iW=A^i,\qquad B^iWVB^j=B^iB^j.
    \]
    Their entrywise images under $f$ are a reduction from $f(B)$ onto
    $f(A)$ on every word, including the empty word.
\end{lemma}
\begin{proof}\leanok
    The homomorphism preserves products, sums and identity matrices.
    Transport the displayed identities and induct on the length of a word,
    inserting $WV$ between its first two letters.
\end{proof}

\begin{lemma}[Local determination of an action scalar]
    \label{lem:mpug_action_scalar_local}
    \lean{MPOTensor.GroupFamily.ActionData.eq_lSymbol_of_isDressedProportional,
        MPSTensor.IsDressedProportional.of_forall_mul,
        MPSTensor.isDressedProportional_complexOfRing_of_forall_mul}
    \leanok
    If two boundaries satisfy $LB^i=zRB^i$ for every letter, then
    $LB^w=zRB^w$ for every nonempty word. The same conclusion follows from
    such identities over a coefficient ring mapped to $\C$.
    If $L,R$ are the two action trees of a normal invariant state of positive
    bond dimension, their scalar is the state's selected $L$-symbol.
\end{lemma}
\begin{proof}\leanok
    Multiply the one-letter identity by the remaining letters. The left
    action tree is a reduction onto the normal state, so it is nonzero
    against arbitrarily long words. This gives uniqueness of the scalar.
\end{proof}

\begin{lemma}[Independence of the relative class from the reduction data]
    \label{lem:mpug_relative_l_gauge}
    \lean{MPOTensor.GroupFamily.ActionData.toBlockActionData,
        MPOTensor.GroupFamily.ActionData.toBlockActionData_lSymbol,
        MPOTensor.GroupFamily.ActionData.lSymbol_eq_blockGauge,
        MPOTensor.GroupFamily.ActionData.relative_lSymbol_isMulCoboundary₂_div}
    \leanok
    Let an exact normal group family fix two normal states of positive bond
    dimensions. For each state choose action reductions, and choose common
    fusion reductions for the group family. The single-state $L$-symbol
    agrees with the one-block $L$-symbol. If the reduction data are changed,
    with common fusion gauge $\beta$ and action gauges $\gamma_A,\gamma_B$,
    then
    \[
        L_A'(g,h)=\frac{\gamma_A(gh)\beta(g,h)}
        {\gamma_A(g)\gamma_A(h)}L_A(g,h).
    \]
    The analogous equation holds for $B$. Thus, writing $Q=L_A/L_B$,
    \[
        Q'(g,h)=\frac{\phi(g)\phi(h)}{\phi(gh)}Q(g,h),
        \qquad \phi(g)=\frac{\gamma_B(g)}{\gamma_A(g)}.
    \]
    In particular $Q'/Q$ is a two-coboundary.
\end{lemma}
\begin{proof}\leanok
    Apply the gauge law for permuted blocks to the one-point block set.
    Its bond identification is the identity. Divide the two gauge equations;
    the common factor $\beta(g,h)$ cancels.
\end{proof}

\begin{definition}[Rational reduction matrices for the combined family]
    \label{def:mpug_controlled_z_rational_reductions}
    \lean{ControlledZGauging.combinedLabelRational,
        ControlledZGauging.combinedRationalTensor,
        ControlledZGauging.stateRational, ControlledZGauging.zeroRational,
        ControlledZGauging.combinedRationalAction,
        ControlledZGauging.combinedRationalProduct,
        ControlledZGauging.combinedStateV, ControlledZGauging.combinedStateW,
        ControlledZGauging.combinedZeroV, ControlledZGauging.combinedZeroW,
        ControlledZGauging.combinedFusionV, ControlledZGauging.combinedFusionW}
    \leanok
    Write $g=(p,q)$, with $p,q\in\{0,1\}$. Over $\Q$, the tensor is
    \[
        T_{p,0}^{ij}=\delta_{ij}(-1)^{pj},\qquad
        (T_{p,1}^{ij})_{lr}=\delta_{ij}\delta_{rj}(-1)^{pj+lj}.
    \]
    The state matrices are the rational forms of $A^i$ and
    $A_0^i=\delta_{i0}$. Their action and stacked product use the same
    Kronecker contractions as the complex tensors. Set
    \[
        V_{p,0}=W_{p,0}=\operatorname{diag}(1,(-1)^p,(-1)^p).
    \]
    For $s=(-1)^p$, set
    \[
        V_{p,1}=\begin{pmatrix}
            0 & 2s & 2s & 0  & 0 & 0 \\
            1 & 0  & 0  & -1 & 0 & 0 \\
            1 & 0  & 0  & 1  & 0 & 0
        \end{pmatrix},\qquad
        W_{p,1}=\frac12\begin{pmatrix}
            0  & 1  & 1 \\
            s  & 0  & 0 \\
            0  & 0  & 0 \\
            0  & -1 & 1 \\
            -s & 0  & 0 \\
            0  & 0  & 0
        \end{pmatrix}.
    \]
    For the product state the caps are $V_{0}=W_{0}=1$ and
    $V_{1}=(1\;0)$, $W_{1}=(1\;1)^T$. Fusion caps are the identity unless
    both controlled-$Z$ bits are one; in that case they are
    $F^<=(1\;0\;0\;0)$ and $F^>=(1\;1\;1\;1)^T$.
\end{definition}

\begin{lemma}[The displayed matrices are actual action and fusion reductions]
    \label{lem:mpug_controlled_z_explicit_reductions}
    \lean{ControlledZGauging.complexOfRing_combinedRationalTensor,
        ControlledZGauging.complexOfRing_stateRational,
        ControlledZGauging.complexOfRing_zeroRational,
        ControlledZGauging.complexOfRing_combinedRationalAction,
        ControlledZGauging.complexOfRing_combinedRationalProduct,
        ControlledZGauging.combinedStateActionData,
        ControlledZGauging.combinedZeroActionData,
        ControlledZGauging.combinedFusionData}
    \leanok
    \uses{def:mpug_controlled_z_rational_reductions, lem:mpug_reduction_coefficient_map}
    The rational tensors map entrywise to the actual combined complex
    tensors and their contractions. The displayed caps give action
    reductions for both states and common fusion reductions for the family.
    Every reduction equation holds on all words.
\end{lemma}
\begin{proof}\leanok
    Check the cap, single-letter and projection-insertion identities over
    $\Q$ for all labels and letters. Apply
    Lemma~\ref{lem:mpug_reduction_coefficient_map}.
\end{proof}

\begin{definition}[The two action trees and their relative factor]
    \label{def:mpug_controlled_z_relative_factor}
    \lean{ControlledZGauging.combinedRationalActTree,
        ControlledZGauging.combinedRationalFuseTree,
        ControlledZGauging.combinedRationalDoubleAction,
        ControlledZGauging.combinedStateLValue,
        ControlledZGauging.combinedRelativeCocycle,
        ControlledZGauging.combinedLClassGauge,
        ControlledZGauging.combinedRelativeLSymbol}
    \leanok
    The rational action tree is $V_g(I\otimes V_h)\alpha^{-1}$ and the
    fusion tree is $V_{gh}(F^<_{g,h}\otimes I)$, with $\alpha$ the canonical
    bond reassociation. They act on $(T_gT_h)\cdot A$.
    For $g=(p,q)$ and $h=(r,s)$ define
    \[
        \lambda(g,h)=(-1)^{ps}4^{qs},\qquad
        \phi(g)=(-1)^{pq}2^q.
    \]
    For arbitrary actual action and common fusion reductions define the
    relative factor $Q(g,h)=L_A(g,h)/L_0(g,h)$.
\end{definition}

\begin{lemma}[The reduction-independent relative action class]
    \label{lem:mpug_controlled_z_relative_l_class}
    \lean{ControlledZGauging.combinedStateActionData_lSymbol,
        ControlledZGauging.combinedZeroActionData_lSymbol,
        ControlledZGauging.combinedRelativeCocycle_eq_lSymbol_div,
        ControlledZGauging.combinedRelativeCocycle_div_cluster_isMulCoboundary,
        ControlledZGauging.combinedRelativeCocycle_isCocycle,
        ControlledZGauging.combinedRelativeCocycle_H2π_eq_cluster,
        ControlledZGauging.combinedRelativeCocycle_H2π_ne_zero,
        ControlledZGauging.combinedRelativeLSymbol_div_computed_isMulCoboundary,
        ControlledZGauging.combinedRelativeLSymbol_isCocycle,
        ControlledZGauging.combinedRelativeLSymbol_H2π_eq_cluster,
        ControlledZGauging.combinedRelativeLSymbol_H2π_ne_zero}
    \leanok
    \uses{lem:mpug_controlled_z_explicit_reductions, lem:mpug_action_scalar_local,
        lem:mpug_relative_l_gauge, def:mpug_controlled_z_relative_factor}
    The explicit reductions give $L_A(g,h)=\lambda(g,h)$ and $L_0(g,h)=1$.
    For the cluster factor system $\omega(g,h)=(-1)^{qr}$,
    \[
        \lambda(g,h)=\frac{\phi(g)\phi(h)}{\phi(gh)}\omega(g,h).
    \]
    Consequently every choice of actual reduction data gives a genuine
    relative two-cocycle $Q$, with
    \[
        [Q]=[\lambda]=[\omega]\ne0
        \quad\hbox{in }H^2(\Z_2\times\Z_2,\C^\times).
    \]
    This identifies the action class with the split-index class. Physical
    gapped-path invariance remains separate; see
    \cite{gap:fbc25_controlled_z_split_index_class_scope}.
\end{lemma}
\begin{proof}\leanok
    Each one-letter comparison of the two action trees has scalar
    $\lambda(g,h)$ for $A$ and scalar one for the product state. Apply
    Lemma~\ref{lem:mpug_action_scalar_local}. The displayed coboundary
    equation is checked on the four group elements. The cluster class is
    nonzero. Lemma~\ref{lem:mpug_relative_l_gauge} then removes all dependence
    on the chosen action and common fusion reductions.
\end{proof}
